\input{../000_header_year.txt}

\begin{document}

%++++++++++++++++++++++++++++++++++++++++++++++++++++++++++++++++++++++++++++++++
%++++++++++++++++++++++++++++++++++++++++++++++++++++++++++++++++++++++++++++++++
%++++++++++++++++++++++++++++++++++++++++++++++++++++++++++++++++++++++++++++++++
\exheader
%++++++++++++++++++++++++++++++++++++++++++++++++++++++++++++++++++++++++++++++++
%++++++++++++++++++++++++++++++++++++++++++++++++++++++++++++++++++++++++++++++++
%++++++++++++++++++++++++++++++++++++++++++++++++++++++++++++++++++++++++++++++++

\noindent
{\bf 極配置（演習）}

次の状態方程式
\begin{align*}
	\frac{d}{dt}
	\left (
	\begin{array}{c}
	 x_1 \\ x_2 
	\end{array}
	\right )
	& = 
	\left (
	\begin{array}{cc}
	 0 & 1 \\
	 2 & -1
	\end{array}
	\right )
	\left (
	\begin{array}{c}
	 x_1 \\ x_2 
	\end{array}
	\right )
	+
	\left (
	\begin{array}{c}
	 0 \\ 1 
	\end{array}
	\right )
	u
\end{align*}
に対して閉ループ系の極が$-4,-5$となるような状態フィードバック
\begin{align*}
	u 
	&=
	\left (
	\begin{array}{cc}
	 f_1 & f_2
	\end{array}
	\right )
	\left (
	\begin{array}{c}
	 x_1 \\ x_2 
	\end{array}
	\right )
\end{align*}
を求めよ．

%---------------------------------------------------------------------------
\newpage

\noindent
{\bf 極配置（演習）}

次の状態方程式
\begin{align*}
	\frac{d}{dt}
	\left (
	\begin{array}{c}
	 x_1 \\ x_2 
	\end{array}
	\right )
	& = 
	\left (
	\begin{array}{cc}
	 0 & 1 \\
	 2 & -1
	\end{array}
	\right )
	\left (
	\begin{array}{c}
	 x_1 \\ x_2 
	\end{array}
	\right )
	+
	\left (
	\begin{array}{c}
	 0 \\ 1 
	\end{array}
	\right )
	u
\end{align*}
に対して閉ループ系の極が$-4,-5$となるような状態フィードバック
\begin{align*}
	u 
	&=
	\left (
	\begin{array}{cc}
	 f_1 & f_2
	\end{array}
	\right )
	\left (
	\begin{array}{c}
	 x_1 \\ x_2 
	\end{array}
	\right )
\end{align*}
を求めよ．
\\

閉ループ系は
\begin{align*}
	\frac{d}{dt}
	\left (
	\begin{array}{c}
	 x_1 \\ x_2 
	\end{array}
	\right )
	& = 
	\left (
	\begin{array}{cc}
	 0 & 1 \\
	 2 & -1
	\end{array}
	\right )
	\left (
	\begin{array}{c}
	 x_1 \\ x_2 
	\end{array}
	\right )
	+
	\left (
	\begin{array}{c}
	 0 \\ 1 
	\end{array}
	\right )
	\left (
	\begin{array}{cc}
	 f_1 & f_2
	\end{array}
	\right )
	\left (
	\begin{array}{c}
	 x_1 \\ x_2 
	\end{array}
	\right )
\\
	&=
	\left (
	\begin{array}{cc}
	 0 & 1 \\
	 2 & -1
	\end{array}
	\right )
	\left (
	\begin{array}{c}
	 x_1 \\ x_2 
	\end{array}
	\right )
	+
	\left (
	\begin{array}{cc}
	 0   & 0   \\
	 f_1 & f_2
	\end{array}
	\right )
	\left (
	\begin{array}{c}
	 x_1 \\ x_2 
	\end{array}
	\right )
\\
	&=
	\left (
	\begin{array}{cc}
	 0 & 1 \\
	 f_1+2 & f_2-1
	\end{array}
	\right )
	\left (
	\begin{array}{c}
	 x_1 \\ x_2 
	\end{array}
	\right )
\end{align*}
特性多項式は
\begin{align*}
	\mbox{det}[sI-(A+BF)]
	& =
	\mbox{det}
	\left (
	\begin{array}{cc}
	 s & - 1 \\
	 - f_1 -2 & s- (f_2-1)
	\end{array}
	\right )
\\
	&= s [s-(f_2-1)] - f_1 -2
\\
	&= s^2 - (f_2-1) s - f_1 -2
\end{align*}
望ましい特性多項式は望ましい極が$-4,-5$であることより
\begin{align*}
	(s+4)(s+5) = s^2 + 9s + 20
\end{align*}
よって係数を比較して
\begin{align*}
	-(f_2-1)=9, \ \ \ -f_1-2 =20
\end{align*}
よって
\begin{align*}
	f_2=-8, \ \ \ f_1=-22
\end{align*}


\end{document}

